\section{When Asking Changes the Answer}\label{app:measurement}
Let $P$ describe latent complete records and let $K_T$ be the distribution
of the recorded transcript given a latent record under form $T$. A frozen
score estimates
\[
 \widetilde\Delta_{q,K}
   =\E_{P,q,K_T}H_T(O_T),
\]
which can differ from the latent $\Delta=\E_P[Yd(X)]$. A branching policy
has its own sequential measurement kernel and hence its own mean score.
Pointwise unbiasedness for unchanged answers does not equate these means
under arbitrary collection-specific kernels. Randomized common anchors
identify mode or position differences for those anchors; identifying the
latent target additionally needs a measurement model or a reference.

We perform an exact semi-synthetic sensitivity calculation on the ten
binary tasks. At query position $j=1,\ldots,6$, independently flip the
answer with probability $\epsilon(j-1)/5$, for
$\epsilon\in\{0,0.01,0.03,0.05,0.1\}$. Labels are unchanged. Fixed forms
ask coordinates in increasing index order. Branching routes on the
\emph{recorded} noisy answers and stops at its leaf; its probability output
and old stump both use those answers. We sum all $2^6$ flip patterns to
obtain exact conditional first and second score moments, averaging over
people and fixed-form assignment. This is an illustrative position model,
not an estimated response-error process for either source survey.

\begin{table*}[t]
\centering
\caption{Position-dependent flips, integrated exactly over every noise pattern. Median absolute latent-target bias and bias divided by the clean branching standard error at $n=1000$. The final-position error rate is $\epsilon$; earlier rates increase linearly from zero.}
\label{tab:measurement}
\small
\setlength{\tabcolsep}{4pt}
\begin{tabular}{lrrrr}
\toprule
Data & $\epsilon$ & Method & Median $|\mathrm{bias}|$ & Bias / clean SE \\
\midrule
SUSHI & 0.00 & Branching & 0.00000 & 0.000 \\
SUSHI & 0.00 & Form scores & 0.00000 & 0.000 \\
SUSHI & 0.03 & Branching & 0.00076 & 0.164 \\
SUSHI & 0.03 & Form scores & 0.00095 & 0.203 \\
SUSHI & 0.10 & Branching & 0.00247 & 0.527 \\
SUSHI & 0.10 & Form scores & 0.00309 & 0.663 \\
SAPA & 0.00 & Branching & 0.00000 & 0.000 \\
SAPA & 0.00 & Form scores & 0.00000 & 0.000 \\
SAPA & 0.03 & Branching & 0.00233 & 0.537 \\
SAPA & 0.03 & Form scores & 0.00187 & 0.455 \\
SAPA & 0.10 & Branching & 0.00777 & 1.787 \\
SAPA & 0.10 & Form scores & 0.00621 & 1.513 \\
\bottomrule
\end{tabular}
\end{table*}


At a last-position flip probability of 0.1, median absolute bias in units
of a clean branching standard error at $n=1000$ is 0.53/0.66 for
branching/form scores on SUSHI and 1.79/1.51 on SAPA. Thus a measurement
effect can dominate the sampling gain, and neither design is uniformly
more robust in this model. Repeating a fixed biased score at larger $n$
reduces its sampling variance but leaves its measurement bias unchanged.

